\sscp{Brief summary of the prehistory of the region}{15-07-04}
Putting it all together,
what does the prehistory of the Wallacean languages in eastern
Indonesia and Timor-Leste look like succinctly? (See \crf{ch:LayAdd} for
details, discussion, and references, along with supporting data
and arguments in chapters 3-14.)

There is evidence in the archaeological record of human habitation
in linguistic Wallacea going back to 42,000 BP for Timor \citep{OConnor2015},
and around 35,000 BP for north Maluku \citep{Bellwood2019a}.
These early travellers were able to make several deep-water crossings
in Wallacea spanning many tens of kilometres.
And they had the (boat?) technology to catch many open ocean
(pelagic) fish that are not easily caught from shore.
This is reflected by thousands of pelagic fish bones from single
archaeological sites on Timor,
along with some of the earliest fish hooks in the world \citep{OConnor2015}.
Some of these early inhabitants stayed around Timor for thousands of years.
We do not know what languages they spoke at the time.
The long time depth,
and the yearly seasonal east
and west prevailing winds means that a scenario restricting human
habitation in the region to a single incursion migrating in a
single direction by a single group is not reasonable.

\largerpage
There is evidence in the archaeological record that several thousand
years before the arrival of the Austronesian-speaking peoples,
a species of cuscus (\it{Phalanger orientalis}) which is native
to the New Guinea region in the east,
was introduced by human transfer into parts of western Maluku
and Timor, where it had not been previously found \citep{FlanneryNitihaminoto1995,Schapper2011}.
There is further evidence that important food crops native to
the New Guinea region,
such as a species of banana,
as well as taro,
sago, breadfruit, and sugarcane were introduced by humans westward
into parts of Maluku
and the Lesser Sundas,
some reaching as far west as western Flores.
Words associated with the banana originating in Papuan languages
in New Guinea related to the pseudo-etymon {\#}muku ‘banana’
are found in three language families stretching from the Bomberai
region (Bird’s Neck) of Papua to western Flores
\citep{DenhamDonohue2009,DonohueDenham2009,DonohueDenham2010,UsherSchapper2022}.
It is likely that these items dispersed through local networks
of contact or trade.

Prior to the arrival of the Austronesian-speaking peoples into
linguistic Wallacea around 4,000 BP,
the languages spoken by inhabitants in the region were predominantly
SOV in their typologies,
with postpositions, preposed possessors,
and post-verbal (or clause-final) negation.
There was a tendency to make an alienable/inalienable distinction
in their nouns, and some made a human/non-human distinction in their third person
pronouns
\citep{Klamer2019,KlamerReesinkVanStaden2008,Reesink2002b,ReesinkDunn2018}.
Evidence shows these non-Austronesian languages were spoken as
far west as eastern Sumbawa \citep{Donohue2007}.
Significant trace evidence is found throughout the region,
even quite far from where Papuan languages are still spoken today,
such as in eastern Flores,
Sabu, Rote, western Timor, Sula,
Buru, Seram, and across southern Maluku.

There is evidence in the, as yet sketchy, archaeological record
that Austronesian-speaking peoples arrived in the Wallacea
region around 4,000 BP.
This is evidenced by certain types of red-slipped pottery
and animal bones associated with the Austronesians.
Specifically it appears that Austronesian-speaking peoples were
in the Timor region by around 3,800 BP \citep{Spriggs2011},
and in north Maluku (Kayoa) around 3,300 BP \citep{Bellwood2019a}.
The 500 year difference is perhaps significant during a period
of rapid Austronesian expansion throughout Indonesia
and the Pacific, and has implications for our understanding of
how they might have spread through Wallacea.
More archaeological work is needed.

The linguistic record indicates the Austronesian languages that
first came into Wallacea were only trivially different from
PMP in their phonologies and lexicons.
In their grammars, they appear to have been mostly SVO,
prepositional, with post-posted possessors
and pre-verbal negation, like many languages currently spoken
in central and western Indonesia.
They did not have a distinction between alienable
and inalienable nouns in their possessive systems,
nor did they make a human/non-human distinction in their third person pronouns.

The linguistic record indicates that Austronesian-speaking peoples
very probably did not enter the region in a single monolithic
advance-by-wave dispersal, but rather leap-frogged into the region,
with different groups starting in different places
and getting off different boats at different island groups at different times.
This is reflected in the discontinuities in the historical phonology,
the morphosyntax, and the lexicon.
In each region, there were different patterns of contact with
non-Austronesian SOV languages, eventually resulting in many different patterns of which features
were adopted where and to what degree.
Over time, the resulting Austronesian languages found in linguistic
Wallacea became very different from Austronesian languages to
the west and north, and reflect a kind of typological hybridization at
the macro level, with many variations of language mixing at the micro level.

The seasonal yearly east
and west prevailing winds,
combined with Austronesian boat technology meant that these newcomers
could travel in any direction at any time,
waves and weather permitting.
There is abundant evidence of many kinds,
of internal migrations going in many directions.
The time-depth of around 4,000 years means that a scenario of
one group arriving one time in one place
and travelling in only one direction is not a reasonable limitation.

The evidence of significant contact-induced change on mostly
SVO languages from typologically SOV languages not only in the
vocabularies, but more significantly in many aspects of the grammar (\srf{02-03})
reflects intense contact and bilingualism over a long period of time \citep{ThomasonKaufman1988}.
The Austronesian-speaking peoples have been shown to be accommodating
of outsiders and willing to adapt
and incorporate other ways of doing things
and talking about things.
So Wallacea seems to have been more of a respectful meeting ground
with inter-marriage and bilingualism, rather than an invasion that pushed out or
killed off the traditional owners.

Different groups in different regions gradually diverged in unique ways.
Along with their different climates,
different flora and fauna, and different patterns of contact
with non-Austronesian languages, they eventually developed into the eight Wallacean subgroups
that we have identified.

The linguistic record seems to indicate that there was a major
incursion by Austronesian-speaking peoples from the west or north
into the Sumbawa-Flores-Lembata region,
becoming what is now the \tsc{Bima-Lembata} subgroup.
Another minor incursion, went from central Sulawesi to Banggai
to the Sula archipelago,
which is now politically part of Maluku.
It is not yet clear where Austronesian-speaking peoples came
into the regions that eventually became the other Wallacean subgroups.

\largerpage
Coming from the west,
the mitochondrial DNA of humans shows a genetic boundary with
a significant increase associated with Melanesian markers around
the Sumbawa-western Flores region.
This increases as one moves eastward.
From eastern Flores, Lembata
and Timor, the mixture with groups showing Asian ancestry is
“almost half-half” \citep{Fricke2019,Klamer2019}.
There is a great deal of genetic mixture in the Wallacean region.
The genetic data are often inconclusive about the directions
of these dispersals \citep{LansingSudoyo2011}.
This is not surprising,
since languages and DNA are independent of each other.

We have not been able to link these diverse Wallacean subgroups
together in any way that supports the idea of a higher-level
“CMP” subgroup, nor even a “CEMP” subgroup.
And we have shown that the arguments that have been proposed
in support of “CEMP” are flawed in a variety of ways,
have been withdrawn, or have been downgraded as due to diffusion,
and therefore are no longer useful for subgrouping.

Yet there are some words that have been identified that are found
in some Wallacean subgroups but not others,
and some are not found west of linguistic Wallacea.
Some of these are also found in Oceanic languages.
Their uneven and non-overlapping distribution does not support
the idea of a common parent language in Wallacea.
But patterns within Wallacea,
for example, such as a number of words found along the southern
islands, but not in geographical central Maluku,
and vice versa, suggest that over the span of 4,000 years there
were regional interlanguages that came
and went, perhaps as their usefulness was superseded by other
languages of wider communication.
We have documented a number of these,
and feel it is reasonable to assume that the uneven distribution
of some words in a limited number of subgroups or only certain
 nodes within several subgroups may be trace evidence of other
interlanguages that have ceased to exist or ceased to function in that capacity.
Vocabulary shared by only some Wallacean subgroups should be
seen as due to regional patterns of diffusion
and contact, rather than as reflecting a parent language that
is common to all Wallacean subgroups.\footnote{
See \citet{Ross2017}
for similar arguments for Oceanic languages.}

We have said that we are unable to link the Wallacean subgroups
together at an identifiable node below PMP.
We have said that the retention of conservative features found
in the region would require reconstructing a proto-language that
is almost identical to PMP,
including trisyllabic roots, word-final consonants,
final vowel-glide sequences, medial CC clusters,
*ə, *q, possibly *h, stative prefix *ma-, causative prefix *pa-,
reciprocal prefix *paR-, applicative *-kV.

However, while the historical phonology
and much of the lexicon of the region are very close to what
we know about PMP,
the grammars of the region do not look like Formosan or Philippine area grammars.
In addition to the typological features associated with SOV Papuan
languages, the languages in the Wallacean subgroups are typologically (mostly) SVO.
Rather than having a voice opposition,
the grammars of most languages pivot around a semantic opposition
of the macroroles of Actor and Undergoer.
In some languages the distinction is marked in the pronominal
systems, and in others it is found in the verbal prefixes.
Split-S languages of various flavours are widespread throughout
the region, but not in all languages.
Verbs of motion and posture (and sometimes experiencer verbs)
are also widespread, which are (optionally) syntactically transitive but semantically
coreferential, covering the semantic range found in active intransitive verbs in Indo-European languages.
From the perspective of the grammar,
it seems unlikely that the Wallacean subgroups link directly to PMP.
It is entirely likely that the various Wallacean subgroups link
back to PMP via different intermediate pathways.
Those pathways are not yet clear.
Once the subgrouping of the western
and northern Austronesian region becomes clearer,
perhaps we will be able to link certain Wallacean subgroups to
other subgroups below PMP through their historical phonology and morphosyntax.
For the moment, we are unable to do that.

Outside our target region,
it appears that there was a group that stayed in one area long
enough to become the ancestral language to what we know as \tsc{POc}.

In the higher-level \it{classification},
if “CMP” does not exist,
and “CEMP” is abandoned,
but EMP is retained,
what are the implications?
And perhaps the idea that SHWNG links closer to \tsc{POc} than
to Wallacean subgroups also needs to be re-examined,
since that issue was proposed in the past based on assumptions
of subgrouping within PMP that have turned out to be without strong foundations.
Or to put it in other words,
what might PMP look like without CEMP?

\begin{figure}[h]\centering
	\includegraphics[width=\columnwidth]{AN-rake-tree}
	\fcp{Rake-like structure of MP and contact zone}{15-13}
\end{figure}

Very briefly, PMP would look like a rake,
just like it does in the Philippines,
Malaysia, and western and central Indonesia
(\citealp{Ross2017}, \citealp[19]{Smith2017b,Smith2024}, \citealp{Klamer2019}).
The difference is that the Wallacean subgroups additionally reflect
a zone of intense but uneven contact with Papuan languages of
a very different typology over a long period of time,
and in a region where they encountered previously unknown flora
and fauna, and new ways of doing
and talking about things.
Keeping in mind that tree diagrams are inadequate to reflect
the complexities we find,
particularly in contact situations,
an attempt is nevertheless made to visualise the situation in
linguistic Wallacea in very broad terms in \frf{15-13}.

How do we account for words
and features in the grammar that are found in some Wallacean
languages and subgroups as well as in Oceanic languages? In the Wallacean
region these do not appear to be attributable to a parent language
common to all Wallacean subgroups.
As we have shown in \srf{02-02}, \srf{15-Features},
\srf{15-04} and in chapters \ref{ch:BL}--\ref{ch:MorLes},
their uneven distribution is attributable to diffusion,
parallel development, and contact,
including local trade networks and regional interlanguages.
Given the evidence of many trade networks
and a high degree of mobility with people moving around the region
in all directions, it is also possible that a few words or features diffused as
a back eddy from the SHWNG region into the areas where other
Wallacean subgroups are spoken.
And some features that begin to appear or circulate in some Wallacean
subgroups seem to have found their way into the language(s) ancestral to \tsc{POc}.

The complex variations, discontinuities, and uneven distribution of numerous features
in the languages of Wallacea show it to be a contact zone between two very different
types of languages, rather than an area that has many daughter languages descended
from a common Austronesian parent language or linkage below PMP.
This zone of contact in the languages mirrors the patterns
Alfred Russel Wallace observed in the fauna of this transition
zone over 150 years ago between two very different biogeographical regions.